\subsection{Sinus und Kosinus am Einheitskreis}
\label{subsec:real_functions_unit_circle}
Der \emph{Einheitskreis} \EnglishTerm{unit circle} hat den Radius $1$.
Wir beginnen bei $(1,0)$ und tragen den Winkel $x$ gegen den Uhrzeigersinn
ab; negative Winkel werden im Uhrzeigersinn gemessen.
\par\noindent\begin{minipage}{\linewidth}
\begin{DefinitionBox}{Sinus und Kosinus für beliebige Winkel}{real_functions_sine_cosine}
  \begin{minipage}[c]{0.49\linewidth}
    Zum Winkel $x\in\R$ gehört auf dem Einheitskreis der Punkt
    \[
      P=(\cos x,\sin x).
    \]
    Der Kosinus ist seine horizontale, der Sinus seine vertikale Koordinate.
    Nach dem Satz des Pythagoras gilt
    \[
      \cos^2x+\sin^2x=1.
    \]
  \end{minipage}\hfill
  \begin{minipage}[c]{0.47\linewidth}
    \centering
    \CachedTikzFigure{content/resources/chapters/real_functions/tikz/functions_unit_circle.tex}
  \end{minipage}
\end{DefinitionBox}
\begin{ExampleBox}[before skip=0.3\baselineskip]{Eine Umdrehung in Viertelschritten}{real_functions_quarter_turns}
  Lesen Sie die Koordinaten an den vier Achsenschnittpunkten des Kreises ab.
  \[
    \begin{array}{c|ccccc}
      x&0&\pi/2&\pi&3\pi/2&2\pi\\ \hline
      \cos x&1&\Cloze{1cm}{0}&\Cloze{1cm}{-1}&\Cloze{1cm}{0}&1\\[0.6em]
      \sin x&0&\Cloze{1cm}{1}&\Cloze{1cm}{0}&\Cloze{1cm}{-1}&0
    \end{array}
  \]
  Nach einer ganzen Umdrehung wird derselbe Punkt wieder erreicht.
\end{ExampleBox}
\end{minipage}

% Handschrift: PDF-Seite 37 (negative Winkel und weitere Vierteldrehung).
Auch negative Winkel und Winkel über einer ganzen Umdrehung folgen
aus dem Einheitskreis. Ergänzen Sie die drei zusätzlichen Spalten:
\[
  \begin{array}{c||c|c|c}
    x&-\pi&-\pi/2&5\pi/2\\\hline
    \sin x&\Cloze{1cm}{0}&\Cloze{1cm}{-1}&\Cloze{1cm}{1}\\[0.6em]
    \cos x&\Cloze{1cm}{-1}&\Cloze{1cm}{0}&\Cloze{1cm}{0}
  \end{array}
\]

\clearpage
\subsection{Periodizität: Derselbe Verlauf wiederholt sich}
\label{subsec:real_functions_periodicity}
Anders als eine Exponentialfunktion behalten Sinus und Kosinus ihre
Richtung nicht dauerhaft bei. Stattdessen wiederholt sich ihr Verlauf.
\begin{FunctionPropertyBox}{Periodizität}{real_functions_periodic}
  Eine Funktion $f:D\to\R$ heißt \emph{periodisch}
  \EnglishTerm{periodic}, wenn es eine Zahl $T>0$ mit $D+T=D$ und
  \[
    f(x+T)=f(x)\qquad\text{für alle }x\in D
  \]
  gibt. Die Zahl $T$ heißt \emph{Periode} \EnglishTerm{period}.
  Die kleinste positive Periode, falls vorhanden, heißt
  \emph{Grundperiode} \EnglishTerm{fundamental period}.
\end{FunctionPropertyBox}
\begin{center}
  \CachedTikzFigure{content/resources/chapters/real_functions/tikz/functions_sine_cosine.tex}
\end{center}
\begin{TheoremBox}{Eigenschaften von Sinus und Kosinus}{real_functions_trig_properties}
  Beide Funktionen sind auf $\R$ stetig, beschränkt und besitzen die
  \textcolor{green!55!black}{Wertemenge} $[-1,1]$ sowie die Grundperiode $2\pi$.
  \[
    \begin{array}{c@{\qquad}c}
      \sin(x+2\pi)=\sin x&\cos(x+2\pi)=\cos x\\[0.6em]
      -1\le\sin x\le1&-1\le\cos x\le1\\[0.6em]
      \sin(-x)=-\sin x&\cos(-x)=\cos x\\
      \text{ungerade}&\text{gerade}
    \end{array}
  \]
  Ihre Nullstellen und ihre Beziehung durch Verschiebung lauten:
  \[
    \begin{aligned}
      \sin x=0&\Longleftrightarrow x=k\pi,\\
      \cos x=0&\Longleftrightarrow x=\frac\pi2+k\pi,
           \qquad k\in\Z,\\
      \cos x&=\sin\left(x+\frac\pi2\right).
    \end{aligned}
  \]
\end{TheoremBox}
Sinus und Kosinus sind auf ganz $\R$ nicht monoton.
Jedoch steigt $\sin x$ auf $[-\pi/2,\pi/2]$ streng monoton;
$\cos x$ fällt auf $[0,\pi]$ streng monoton.
Diese Intervalle werden für die Umkehrfunktionen wichtig.

% Handschrift: PDF-Seite 39 (beide Verschiebungsbeziehungen und Vorzeichenwechsel).
\begin{KeyPointBox}{Verschiebungen verbinden die Winkelfunktionen}{real_functions_trig_shifts}
  Eine Vierteldrehung verschiebt Sinus und Kosinus ineinander;
  eine halbe Drehung kehrt ihr Vorzeichen um:
  \[
    \begin{aligned}
      \cos x&=\sin(x+\pi/2),&\sin x&=\cos(x-\pi/2),\\[0.6em]
      \sin(x\pm\pi)&=-\sin x,&\cos(x\pm\pi)&=-\cos x.
    \end{aligned}
  \]
  \begin{center}
    \CachedTikzFigure{content/resources/chapters/real_functions/tikz/functions_trig_shifts.tex}
  \end{center}
\end{KeyPointBox}
